\section{Principalization on the original global Nadel scheme}

\begin{lemma}\label{lem:kn-nadel-generic-nonzero-everywhere}
\lean{EnvelopeProperty.Vanishing.NadelPrincipalization.nonzeroEverywhere_of_generic}
\leanok
On an integral scheme $X$, if the original ideal $I$ is nonzero at
the generic point, it is nonzero at every stalk.
\end{lemma}
\begin{proof}
The generic point specializes to each $x$.  The generic stalk ideal is
the extension of $I_x$ along the native specialization map.  If $I_x$
were zero, its extension would be zero, contradicting the original
generic nonvanishing.  This uses the original ideal and its actual stalk
maps and introduces no resolution datum.
\end{proof}

\begin{definition}\label{def:kn-nadel-original-global-bp}
\lean{EnvelopeProperty.Vanishing.NadelPrincipalization.originalTriple,EnvelopeProperty.Vanishing.NadelPrincipalization.sequence}
\leanok
\uses{lem:kn-nadel-generic-nonzero-everywhere}
For the original integral smooth projective $f:X\to\operatorname{Spec}k$
of dimension $n$ and original ideal $I$ nonzero at the generic point,
construct the algebraic principalization triple with precisely this $X$,
$f$, $I$, $n$ and the empty boundary.  Denote its actual Hironaka BP
sequence by $S$.
\end{definition}
\begin{proof}
Original projectivity gives finite type, quasi-compactness and
separatedness.  Original smoothness gives the prescribed pure relative
dimension and the empty SNC boundary.  The lemma proves the required
nonzero-at-every-stalk condition internally.  Apply the existing proved
algebraic BP construction to this actual triple.
\end{proof}

\begin{theorem}\label{thm:kn-nadel-original-global-principalization}
\lean{EnvelopeProperty.Vanishing.NadelPrincipalization.sequence_isSmooth,EnvelopeProperty.Vanishing.NadelPrincipalization.stage_smoothOfRelativeDimension,EnvelopeProperty.Vanishing.NadelPrincipalization.totalTransform_isIdealOfSncDivisor,EnvelopeProperty.Vanishing.NadelPrincipalization.last_isProjectiveFactorization,EnvelopeProperty.Vanishing.NadelPrincipalization.isIso_restrict_compl}
\leanok
\uses{def:kn-nadel-original-global-bp}
The produced finite sequence consists of smooth-centre blowups, its
last structure morphism $S.\mathrm{composite}\,f$ is smooth of the same
dimension $n$ and has an actual finite projective embedding, and the
total transform $I\mathcal O_{S.\mathrm{last}}$ is the ideal of an SNC
divisor.  Its composite is an isomorphism over the original complement
of the support of $I$.
\end{theorem}
\begin{proof}
The proved centre theorem of algebraic principalization gives smoothness
of every actual centre for this triple.  The smooth-centre blowup
dimension theorem propagates the original $n$ through the finite sequence.
The original projective embedding and the proved polynomial projectivity
of a blowup succession construct a projective factorization of the last
structure morphism.  The BP monomial theorem applies to the original
ideal, giving the actual SNC total transform.  The BP isomorphism-locus
theorem gives the complement of the union of the original ideal support
and the singular locus of the original boundary; the empty boundary has
empty singular locus, so this is precisely the original ideal complement.
\end{proof}

\begin{lemma}\label{lem:kn-nadel-original-principalization-integral}
\lean{EnvelopeProperty.Vanishing.NadelPrincipalization.genericPoint_mem_unitLocus,EnvelopeProperty.Vanishing.NadelPrincipalization.unitLocus_dense,EnvelopeProperty.Vanishing.NadelPrincipalization.last_nonempty,EnvelopeProperty.Vanishing.NadelPrincipalization.last_isIntegral,EnvelopeProperty.Vanishing.NadelPrincipalization.composite_isProper,EnvelopeProperty.Vanishing.NadelPrincipalization.exists_dense_isIso}
\leanok
\uses{thm:kn-nadel-original-global-principalization}
The original generic point belongs to the actual unit locus of $I$.
This open is dense, the actual last scheme is nonempty and integral,
and the actual composite $S.\mathrm{last}\to X$ is proper and isomorphic
over that same dense open.
\end{lemma}
\begin{proof}
The generic stalk is the original function field, so its maximal ideal
is zero.  Generic nonvanishing of $I$ therefore excludes the generic
point from its support.  A nonempty open in the original integral scheme
is dense.  Its inverse under the proved actual isomorphism supplies a
point on the last scheme.  For a finite blowup sequence on an integral
scheme, final nonemptiness forces each intermediate blowup to be
nonempty, hence its centre has proper support; the proved native blowup
integrality theorem then propagates integrality.  Properness follows
from the actual finite succession of blowups on the original Noetherian
scheme.  The original support complement is the asserted dense
isomorphism locus.
\end{proof}

This constructs global principalization from the original geometric data.
It does not assert the global multiplier pushforward formula, birational
canonical acyclicity, or higher-direct-image local vanishing.  Those remain
the separate obligations of Lemma~\ref{lem:kn-nadel-original-resolution}
and Lemma~\ref{lem:kn-nadel-local-vanishing}.
